\documentclass[../../main.tex]{subfiles}
\begin{document}

\paragraph{【新課程】}
箱の中に，1から9までの番号をひとつずつ書いた9枚のカードがある．
それらをよくまぜて，その中から1枚ずつ続けて全部を取り出し，
取り出した順に新しく1から9までの番号をつける．
このとき，新しくつけられる番号が前もってつけられている番号に一致するカードが，
ちょうど5枚できる確率を求めよ．

\paragraph{【旧課程】}
点$A_0$を一端とする半直線$l$上に点$A_1$，$A_2$が与えられていて，
$A_0A_1=1$，$A_0A_2=8$であるとする．
いまこれから，点$A_n$ ($n=3, 4, 5, \cdots$)を次の手順で定める．
\begin{enumerate}[(1)]
  \item $A_{n-2}$，$A_{n-1}$のうち$A_0$に近い方を$X$，遠い方を$Y$とする．
    線分$A_0Y$を直径とする半円と，$X$における$l$の垂線との交点をとってこれを$Z$とする．
  \item 次に，$A_0$を中心とする半径$A_0Z$の円と，
    半直線$l$との交点をとってこれを$A_n$とする．
\end{enumerate}
このとき，点$A_n$はある定点$A$に限りなく近づく．
線分$A_0A$の長さを求めよ．

\end{document}
